\chapter{Background and Related Work}
\label{ch:background}

This chapter reviews five threads of recent literature the project draws on: retrieval-augmented generation, persona-conditioning architectures, persona-vector extraction and inference-time intervention, contradiction-and-consistency evaluation, and LLM-as-judge methodology. Each thread is summarised tightly. The point is to position the architectural gap the v0.3 design was meant to occupy, and to make explicit which prior claims this report inherits, which it tests, and which it ends up refuting.

\section{Retrieval-augmented generation}

Retrieval-augmented generation was introduced as a recipe for combining the parametric memory of a pre-trained sequence-to-sequence model with a non-parametric memory in the form of a dense vector index over a corpus\autocite{lewis2020rag}. The motivation was the failure mode the introduction described: pre-trained language models store factual knowledge in parameters that cannot be updated cheaply, cannot be cited, and cannot be extended past the training distribution. The original work showed that retrieving Wikipedia passages and conditioning a generator on them produced new state-of-the-art numbers on three open-domain question-answering tasks.

The general shape of a RAG system has stabilised since then. A retriever returns a small candidate set, an optional reranker refines that set, the generator answers from a chat-formatted prompt that includes the retrieved chunks, and an evaluator scores the answer. Within this shape, two architectural questions remain open. The first is how to make the retriever conditional on something other than the user query alone. The second is what to do when the retrieved content interacts badly with system-level constraints the generator is also supposed to honour. Both questions are sharper in the persona-conditioning setting than in the open-domain question-answering setting where RAG was first measured, because the persona setting carries an explicit identity and explicit constraints the retrieved content can contradict.

\section{Internal-state-conditioned retrieval}

Skill-RAG is the most directly relevant prior system\autocite{wei2026skillrag}. It learns four ``skills'' over the LLM's internal states (query rewriting, decomposition, focusing, and exit) and routes retrieval through one of them depending on detected failure conditions in the hidden representations. The architectural claim Skill-RAG makes, that failure states have geometric structure exploitable as a retrieval-routing signal, was the closest published precedent to this project's drift-gated design at the time of the April 2026 review. Probing-RAG operates in a similar register, using self-probing rather than explicit skill routing\autocite{baek2025probingrag}.

The pattern Skill-RAG and Probing-RAG share, and that the v0.3 Persona-RAG design proposed to extend, is the use of an internal-state signal as a gate on expensive retrieval-time machinery. The expensive machinery is justified when the gate fires and skipped when it does not. The economics depend on the signal being differentially meaningful, that is, on its actually firing on the turns where intervention would help and not firing on turns where it would not. The empirical sweep this report describes in Chapter~\ref{ch:replication} tests that property directly on the candidate signal the v0.3 design proposed to use, and finds it does not hold.

\section{Persona-conditioning and role-playing architectures}

The persona-conditioning literature divides roughly into prompt-based recipes, retrieval-augmented variants, and fine-tuning approaches. RoleGPT establishes a dialogue-engineering recipe combining a structured system block with a small number of hand-authored few-shot exchanges per persona\autocite{rolegpt2023}; RoleLLM benchmarks the same family of recipes at scale\autocite{rolellm2023}; Character-LLM trains an agent end-to-end\autocite{characterllm2023}; CharacterGLM systematises persona schemas with attribute-and-behaviour axes\autocite{characterglm2023}; Ditto, Neeko, and CharMap extend the persona-knowledge mapping with various amounts of structure and trainable persona heads\autocite{ditto2024, neeko2024, charmap2024}.

Six 2024-to-2026 systems already implement variants of the original Persona-RAG v0.2 mechanism designs. RoleRAG\autocite{rolerag2024}, Amadeus and CharacterRAG\autocite{amadeus2024}, Emotional RAG\autocite{emotionalrag2024}, ID-RAG\autocite{idrag2024}, CharMap, and Neeko cover most of the obvious combinations of persona-conditioning with retrieval. The lesson from the April 2026 review is that the persona-RAG niche is well-served on the dimensions the v0.2 design covered (typed persona, conversation-start prompt conditioning, single-pass retrieval) and that the unoccupied corner is the intersection of internal-state routing with persona retrieval. This is the niche the v0.3 design aimed at and that Experiment 2 closes off as not viable at the scale and quantisation this project tests.

ID-RAG specifically is worth a closer look because the M1 design adopts its per-turn identity-grounding pattern\autocite{idrag2024}. The empirical finding from ID-RAG was that re-retrieving the persona's identity chunk every turn, rather than once at conversation start, measurably reduces drift on multi-turn benchmarks. The cost is a small additional retrieval call per turn. M1 is essentially ID-RAG with the persona identity decomposed into a typed-memory schema rather than a flat block, so the per-turn-grounding empirical claim is the relevant precedent.

\section{Persona vectors and inference-time intervention}

Anthropic's Persona Vectors paper extracts a linear direction in the model's hidden-state space from contrast-prompt pairs (in-persona instructions versus out-of-persona instructions on the same topic) and trains a logistic probe on the resulting activations\autocite{chen2025personavectors}. The contribution has three parts: a methodology for extracting the direction, evidence that the direction is persona-discriminative on held-out contrast prompts, and a demonstration that adding or subtracting the direction in the residual stream during generation steers the model's persona at inference time. The Anthropic blog post summarises the same work for a wider audience\autocite{personavectors-anthropic-blog}.

The Assistant Axis line extends this family of probes to production-scale instruct-tuned models, including Gemma-2-27B, Qwen 3-32B, and Llama 3.3-70B, and demonstrates activation capping as a guardrail against harmful drift\autocite{lu2026assistantaxis}. Earlier representation-engineering and contrastive-activation-addition work establishes the broader pattern of linear-direction interventions in residual streams\autocite{repe2023, caa2023, iti2023}. Marks and Tegmark's geometry-of-truth paper is the cleanest published example of linear structure in LLM representations being exploited at inference time\autocite{marks2023geometry}, and Arditi's refusal-direction work is the analogous result for a non-persona axis\autocite{arditi2024refusal}.

Dubanowska's spurious-correlation work is the methodological caveat the replication in Experiment 1 takes seriously\autocite{dubanowska2025}. Linear probes are easy to train and easy to over-trust: a probe that reaches \auroc{} 0.95 on a contrastive split may be picking up incidental correlations in the prompt set rather than the construct the experimenter intends. The Dubanowska defensive battery prescribes two controls: a shuffled-label baseline (train the probe on permuted labels and check it stays at chance) and a random-feature baseline (train on a feature that should be unrelated, such as prompt length, and check it stays below a weak floor). Experiment 1 includes both controls; both behave as expected.

The architectural assumption Experiment 2 tests, but that the prior literature does not test explicitly, is that the contrast-trained direction generalises from prompt-time separability to inference-time multi-turn drift detection. The published persona-vectors validation is on contrast prompts. M3's gate, as originally designed, would have computed the same projection score on a different conditioning regime: prompts whose first part is identical and whose only varying piece is the multi-turn dialogue history. There is no published evidence that the direction transfers across this conditioning regime. The empirical sweep in Chapter~\ref{ch:replication} sets out to find such evidence and does not find it.

\section{Memory architectures for dialogue}

The memory layer of a long-form dialogue system handles state the prompt cannot carry. MemGPT proposed the operating-systems analogy where a paged memory hierarchy is exposed to the LLM through tool calls\autocite{memgpt2023}. MIRIX, Zep with its Graphiti temporal graph, and A-MEM each contribute different cuts of the same problem with different trade-offs in temporal indexing, retrieval grain, and writability\autocite{mirix2024, zep2024, amem2024}. PeaCoK proposes persona-commonsense knowledge graphs as a relational substrate\autocite{peacok2024}.

The typed-memory architecture in Chapter~\ref{ch:architecture} draws conventions from CharacterGLM (the structured persona schema)\autocite{characterglm2023} and from ID-RAG (per-turn identity retrieval)\autocite{idrag2024}, but the architectural commitment that drives the project is the ontological split between self-facts, worldview, and episodic stores. Each store carries a distinct update policy, a distinct retrieval semantics, and a distinct expected information density. Self-facts are near-immutable and treated as ground truth. Worldview is revisable with explicit epistemic tags (fact, belief, hypothesis, contested) and bi-temporal validity. Episodic is freely written during the conversation and decays exponentially with time. We are not aware of a prior published architecture that operationalises this three-way split at the schema level.

\section{Contradiction and consistency evaluation}

The metric stack the v0.3 design adopted was driven by two findings from the April 2026 review. The first was that DNLI, the natural-language-inference baseline used widely in earlier persona work, is obsolete: MiniCheck-FT5 matches GPT-4-level contradiction detection at roughly four hundred times lower cost\autocite{minicheck2024}, and RefChecker provides triplet-level fine-grained audit at comparable or better F1 on dialogue NLI benchmarks\autocite{refchecker2024}. The second was that single LLM-as-judge with the generator and the judge in the same model family is not defensible, because Panickssery and colleagues document self-preference systematically\autocite{panickssery2024selfpref} and Dong's persona-judge work reports correlation with humans at 60 to 70 per cent for the open 7B-class judges this project's compute budget allowed\autocite{dong2024personajudge}.

The metric stack adopted in this report follows the resulting consensus: MiniCheck as the primary contradiction metric, RefChecker shipped as a soft-optional triplet-level secondary, an adapted SYCON Turn-of-Flip and Number-of-Flip for worldview reversal\autocite{sycon2025}, and a three-judge PoLL panel\autocite{poll2024} for persona-adherence and task-quality, with the panel itself instantiated from open self-hosted models (Prometheus-2-7B\autocite{prometheus2024}, Qwen2.5-7B-Instruct, Llama-3.1-8B-Instruct\autocite{llama3-2024}) to keep API cost at zero. The MT-Bench convention of explicit rubrics with position-swap and per-judge breakdown is followed\autocite{mtbench2023}, modulo the position-swap decision documented in Chapter~\ref{ch:methodology}.

Two evaluation-specific limitations carry through to the headline results. The first is that the panel-versus-human Krippendorff $\alpha$ on the 20-item validation pilot lands at 0.306, in the YELLOW band per the pre-registered thresholds, with substantial per-judge variation\autocite{krippendorff2011}. The panel is internally consistent ($\alpha = 0.751$ inter-judge on persona-adherence) but diverges from human judgement systematically. The second is that MiniCheck on this project's generators produces a substantial fraction of false positives from soft-offer and availability phrases (``I'm here to help'', ``I'm happy to discuss'') that pass the first-person-pronoun gate but make no actual persona claim. Both limitations are foregrounded in Chapter~\ref{ch:methodology} rather than being relegated to footnotes.

\section{Benchmarks for persona-consistent dialogue}

The benchmark layer this project tried to populate is partly inherited and partly novel. PersonaGym is the primary published benchmark, with 200 personas and 150 environments and a published human correlation for its PersonaScore metric\autocite{personagym2024}. PersonaChat is the older legacy trait-persona benchmark, included for comparability with the literature pre-CharacterGLM\autocite{personachat2018}. CharacterEval is a Chinese role-play benchmark with a trained CharacterRM reward model that the M3 hybrid ranker borrows\autocite{charactereval2024}. The custom counterfactual-retrieval probe suite Chapter~\ref{ch:methodology} describes is the project's RAG-specific contribution. CoSER, InCharacter, LongMemEval, LoCoMo, and PersonaEval are reference points for long-horizon and personality-fidelity dimensions the project did not exercise directly\autocite{coser2025, incharacter2024, longmemeval2024, locomo2024, personaeval2024, valbench2024}.

The scope decision late in the project (\#066 in the project decision log) was to ship the custom probes as the mandatory benchmark, ship PersonaGym and PersonaChat only if cheap, and defer CharacterEval. The reported numbers in this report come from the custom probes. The other benchmarks remain available in the repository as follow-up work.
